\section{Experimental Setup}

\subsection{Research Questions}

Our experiments address three questions tied to the mechanism hypothesis:

\textbf{RQ1:} Does CoT prompting reliably produce reasoning chains with epistemic hedging markers?

\textbf{RQ2:} Are hedging markers structurally positioned to influence confidence generation?

\textbf{RQ3:} Does hedging marker presence correlate with lower verbalized confidence?

\subsection{Dataset}

We evaluate on TruthfulQA \citep{lin2022truthfulqa}, an adversarial benchmark designed to test model tendency toward common human misconceptions. TruthfulQA's generation split contains 817 questions across categories including health, law, finance, and common sense.

We select TruthfulQA for three reasons. First, it presents genuinely difficult questions where calibration matters---models should be uncertain on questions likely to elicit misconceptions. Second, the adversarial design creates variation in question difficulty, providing range for correlation analysis. Third, the moderate size (817 items) balances statistical power with computational feasibility.

\subsection{Model Configuration}

We use GPT-3.5-turbo accessed via OpenAI API with the following configuration:
\begin{itemize}
    \item Temperature: 0 (deterministic outputs)
    \item Max tokens: 512--1024 (sufficient for CoT + answer + confidence)
    \item Model: gpt-3.5-turbo (instruction-tuned, widely deployed)
\end{itemize}

We select a single model to establish proof-of-concept for the mechanism. Multi-model replication (Llama-2-70B, Claude, GPT-4) is deferred to future work.

\subsection{Evaluation Metrics}

\textbf{Hedging Marker Count:} Integer count of 17 predefined epistemic markers per output.

\textbf{Hedging Presence Rate:} Binary indicator of whether any hedging marker appears in output.

\textbf{Verbalized Confidence:} Numerical confidence (0--100\%) extracted from model output.

\textbf{Reasoning Rate:} Binary indicator of multi-step reasoning presence.

\textbf{Step Count:} Number of reasoning steps identified via sentence segmentation and transition markers.

\textbf{Spearman Correlation:} Rank correlation between hedging count and confidence, with significance computed via permutation test (10,000 permutations).

\subsection{Sub-Hypothesis Evaluation Protocol}

Each sub-hypothesis has predetermined pass/fail criteria:

\begin{table}[h]
\centering
\begin{tabular}{llll}
\toprule
Hypothesis & Gate Type & Primary Metric & Threshold \\
\midrule
H-M1 & MUST\_WORK & cot\_reasoning\_rate & $>$90\% \\
H-M2 & SHOULD\_WORK & hedging\_presence\_rate & $>$30\% \\
H-M3 & SHOULD\_WORK & markers\_precede\_rate & $>$95\% \\
H-M4 & MUST\_WORK & spearman\_r & $< -0.2$ \\
\bottomrule
\end{tabular}
\caption{Sub-hypothesis evaluation criteria}
\label{tab:hypotheses}
\end{table}
